% !TeX root = ../main.tex

\chapter{Evaluation}
\label{ch:evaluation}

This chapter evaluates the tiered compilation architecture. Section~\ref{sec:eval:method} introduces the benchmark suite, describes the strengthening procedure that produces the suite used for measurement, defines the measurement protocol, and lists the hardware and parameters used. Section~\ref{sec:eval:workloads} groups the kernels by execution shape, so that each tier-2 mechanism's effect can be read off the group it targets. Section~\ref{sec:eval:llvm} establishes the whole-module LLVM JIT baseline as the upper bound on steady-state throughput. Sections~\ref{sec:eval:step1}--\ref{sec:eval:step3} measure the tiered architecture in three configurations that correspond to the three design steps. Section~\ref{sec:eval:tt} reports cold-start latency and end-to-end wall-clock time.

\section{Methodology}
\label{sec:eval:method}

\subsection{The sightglass benchmark suite}
\label{sec:eval:method:sg}

Sightglass is the WebAssembly benchmark suite maintained by the Bytecode Alliance. The suite ships a collection of kernels drawn from cryptographic primitives, parsing and serialization workloads, image and string processing, hashing, and small synthetic compute loops. Each kernel is compiled to WebAssembly through the standard toolchain and packaged with an input file and a golden output. The kernels are diverse enough to expose backend code-quality differences, and small enough that each kernel completes in seconds rather than minutes.

The kernels are designed for cheap regression-testing. Most upstream kernels finish in well under a second on a modern host. The brevity is appropriate for correctness regression. The brevity is a measurement hazard for tier-2 evaluation. Tier-2 works by executing under the tier-1 baseline while a worker thread compiles the LLVM-grade replacement, then switching to the optimized code once it is ready. If the entire kernel completes before the switch, the measurement reflects only tier-1 work and reveals nothing about tier-2 code quality. Even when the switch fires, a residual tier-1 prefix of comparable size to the post-switch tail confounds the steady-state signal.

\subsection{The sightglass-strong suite}
\label{sec:eval:method:strong}

To give tier-2 a measurable runway, this thesis strengthens every kernel so that its tier-1-only run-time lands in a 5- to 10-second band, with a target near 8 seconds. The strengthened suite is referred to as \emph{sightglass-strong}. A few seconds of post-switch execution is enough to amortize the worker's compile cost and to distinguish real steady-state speedup from measurement noise. A tight band keeps overall CI time bounded.

Three strengthening patterns cover most kernels. The first pattern increases an internal iteration constant defined at the top of the kernel source, used for shootout-style kernels whose inner loop is bounded by a single macro. The second pattern bumps the literal in the kernel's outer benchmark loop, used for kernels whose driver already loops over the work region. The third pattern wraps the benchmark region in a fresh outer loop, used for Rust and C kernels that originally measured a single execution. A small number of kernels require additional work: out-of-bounds fixes uncovered by the larger iteration count, optimization barriers to prevent the compiler from hoisting the kernel out of the surrounding loop, and replacement of nondeterministic output (such as wall-clock prints) so that the golden output remains stable. The per-kernel strengthening table is reproduced in Appendix~B.

The thirty-three kernels in sightglass-strong cover the same names as the upstream suite, except for two large kernels (spidermonkey and tinygo) that are excluded. The kernel \texttt{noop} is intentionally retained at its upstream size as a harness-overhead reference; its run-time is on the order of one microsecond.

Goldens for sightglass-strong were regenerated by running the strengthened kernels under an independent runtime---wasmtime 43.0.1---rather than under WasmEdge itself. The independent oracle means that any regression in WasmEdge's parser, baseline JIT, or tier-2 pipeline surfaces as a golden mismatch rather than silently corrupting both the runtime output and the reference output.

\subsection{Measurement protocol}
\label{sec:eval:method:protocol}

Each measurement reports a three-run median to reduce single-run variance. All measurements are taken on the same host with the runtime built in release-optimized configuration and the tier-1 SSA backend at its standard optimization level. Three metrics are reported. \emph{WorkTime} (\texttt{WT}) is the harness-measured execution time of the benchmark region. \emph{Instantiation latency} (\texttt{Inst.Lat}) is the upfront compile and bind cost paid before the first wasm instruction runs. \emph{Time-to-value} (\texttt{TtV}) is the end-to-end wall-clock time (\texttt{Inst.Lat + WT}).

All ratios in this chapter follow the convention that values greater than one indicate a tier-2 win. The two ratios cited most often are \texttt{t1/t2} (tier-1 work-time divided by tier-2 work-time) and \texttt{LLVM/t2} (whole-module LLVM JIT work-time divided by tier-2 work-time). The configuration of interest is the second. The ratio \texttt{LLVM/t2} measures how close the tiered architecture comes to delivering whole-module LLVM JIT's steady-state throughput.

\subsection{Experiment setup}
\label{sec:eval:method:setup}

All measurements were taken on a single workstation-class host with no other interactive workload running. The host has a 13th-generation Intel Core i9-13900K processor (24 physical cores spanning performance and efficiency partitions, 32 hardware threads, 5.8\,GHz turbo) and 128\,GiB of DDR5 memory, running Ubuntu 22.04 LTS on Linux kernel 6.8. Each sightglass-strong run executes single-threaded; the worker thread on which tier-2 compilation runs is the only other thread of interest, and it is co-located with the user thread on the same host but never contends for the same core during the benchmark region.

WasmEdge is built in release-optimized configuration. The tier-1 SSA backend runs at its standard optimization level, which is the level the rest of the IR JIT pipeline uses. The whole-module LLVM JIT baseline runs through the unmodified WasmEdge LLVM frontend at its default optimization level, which is the level used for the production AOT path.

Three configurations are measured. Each configuration shares the same tier-1 baseline; the configurations differ only in which tier-2 mechanisms are enabled. Table~\ref{tbl:eval:params} summarizes the parameters.

\begin{table}[h]
\centering
\small
\begin{tabular}{l p{3.5cm} p{3.5cm} p{3.0cm}}
\toprule
Configuration & Tier-2 promotion & Function-entry threshold & OSR threshold \\
\midrule
Tier-1 only        & disabled                    & ---           & --- \\
Step 1 (hot callees) & enabled, leaf + direct callees & 10 calls   & legacy per-loop counter \\
Step 2 (walk-up + BFS) & enabled, walk-up + BFS  & 10 calls      & OSR disabled \\
Step 3 (+ OSR)        & enabled, walk-up + BFS  & 10 calls      & 5000 iterations \\
LLVM JIT baseline  & not applicable              & ---           & --- \\
\bottomrule
\end{tabular}
\caption{Configurations used in this evaluation. The function-entry threshold is the number of calls a tier-1 function counts before triggering tier-2 promotion. The OSR threshold is the number of back-edge iterations a tier-1 loop counts before triggering an OSR continuation compile.}
\label{tbl:eval:params}
\end{table}

The two thresholds are kept fixed across kernels. Step 1 corresponds to the earliest tier-2 design measured on commit \texttt{89e83770}; Steps 2 and 3 correspond to the current head commit \texttt{2dc9ef7b}. The tier-1 baseline numbers used in every step come from the current head, so the step-to-step deltas reflect changes to the tier-2 pipeline rather than incidental changes to tier-1.

The two threshold values are calibrated to the call-count and iteration-count distributions of the strengthened suite. The strengthening procedure tunes each kernel to a 5--10 second tier-1 work-region, which forces every kernel's hot functions to execute many thousands to many millions of times during measurement. A function-entry threshold of 10 therefore saturates within milliseconds of kernel start on every kernel, while remaining well above the call counts of one-time initialization, validation, and cleanup helpers that should not be queued. The OSR threshold sits an order of magnitude below the shortest one-shot-caller loop in the suite: the relevant iteration counts span 60{,}000 (\texttt{shootout-ctype}, the shortest) through 1.2 million (\texttt{shootout-base64}), 35 million (\texttt{shootout-matrix}), 50 million (\texttt{shootout-ratelimit}), and two billion (\texttt{shootout-random}). A threshold of 5000 fires on every one of them---leaving at most 8\% of the loop in tier-1 on \texttt{shootout-ctype} and less than one part in ten thousand on \texttt{shootout-random}---while remaining above the iteration counts of short bookkeeping loops that would otherwise fill the worker queue with continuations the user never re-enters. Both choices are robust over roughly a decade in either direction; the per-kernel statistics are several orders of magnitude away from the threshold on both sides.

\subsection{Excluded kernels}
\label{sec:eval:method:excluded}

The kernel \texttt{noop} is excluded from every aggregate. Its harness-overhead microsecond cannot be averaged meaningfully with seconds-scale kernels. The kernel \texttt{shootout-ackermann} is bimodal on the test host: depending on dynamic factors that have not been fully isolated, it lands in one of two well-separated time bands per run. The bimodality dominates per-kernel ratios involving \texttt{shootout-ackermann}, so this kernel is excluded from the geometric means reported in this chapter. The remaining thirty-one kernels constitute the analysis set.

\section{Workload Characterization}
\label{sec:eval:workloads}

The thirty-one analysis kernels divide into four execution-shape categories. The categorization clarifies which tier-2 mechanism each kernel responds to. The categories are not strictly disjoint: two kernels (\texttt{shootout-ctype} and \texttt{blind-sig}) sit in two categories simultaneously, and the design's full effect on those kernels requires both mechanisms to engage.

\paragraph{Frequent-call workloads.} The kernel invokes a hot function many times, with each call short relative to the total run-time. The function's entry counter saturates quickly, the entry-table swap captures every subsequent call, and the tier-1 prefix preceding the swap is negligible relative to the steady-state tail. Members: \texttt{blake3-scalar}, \texttt{bz2}, \texttt{gcc-loops}, \texttt{hashset}, \texttt{pulldown-cmark}, \texttt{quicksort}, \texttt{regex}, \texttt{shootout-gimli}, \texttt{shootout-heapsort}, \texttt{shootout-keccak}, \texttt{shootout-memmove}, \texttt{shootout-seqhash}, \texttt{shootout-sieve}, \texttt{shootout-switch}, \texttt{shootout-xblabla20}, \texttt{shootout-xchacha20}.

\paragraph{Hot-leaf workloads.} The threshold-tripping function sits several layers below the actual hot work. A small leaf helper crosses its counter first, but the surrounding caller---which contains the bulk of the run-time---is what tier-2 needs to compile to deliver speedup. Promoting only the leaf yields a tier-2 body that calls back into tier-1 for every helper outside the batch, blocking inlining and pinning the call sites at bridge-thunk cost. Members: \texttt{blind-sig}, \texttt{rust-compression}, \texttt{rust-json}, \texttt{rust-protobuf}, \texttt{shootout-ctype}, \texttt{shootout-ed25519}.

\paragraph{One-shot caller workloads.} The kernel invokes its outer function exactly once and spends nearly all of its run-time inside one long-running loop. Function-entry promotion delivers no speedup on these kernels, because the function never returns to its prologue and never re-enters through the swapped slot. Members: \texttt{blind-sig}, \texttt{shootout-base64}, \texttt{shootout-ctype}, \texttt{shootout-matrix}, \texttt{shootout-minicsv}, \texttt{shootout-random}, \texttt{shootout-ratelimit}.

\paragraph{Resistant workloads.} The kernel has either too little total run-time for the worker to amortize its compile cost, or structural properties that block optimizer wins. Members: \texttt{richards} (short total run-time), \texttt{rust-html-rewriter} (short total run-time), \texttt{shootout-fib2} (pure recursive call), \texttt{shootout-nestedloop} (innermost loop with no inline opportunity).

\section{LLVM JIT Baseline}
\label{sec:eval:llvm}

Whole-module LLVM JIT establishes the upper bound on steady-state throughput. Its work-time numbers are the divisor for every \texttt{LLVM/t2} ratio reported in this chapter. The baseline also defines the cost structure that the tiered architecture is designed to avoid.

LLVM JIT compiles every defined function before any wasm instruction executes. Its instantiation latency therefore scales with module size and the complexity of individual function bodies. Three groups appear in the LLVM JIT instantiation-latency distribution. Small modules with few or simple functions instantiate in tens of milliseconds (\texttt{shootout-gimli} 14 ms, \texttt{shootout-heapsort} 87 ms). Mid-sized modules instantiate in hundreds of milliseconds (\texttt{shootout-keccak} 368 ms, \texttt{shootout-matrix} 289 ms, \texttt{pulldown-cmark} 2108 ms). Large modules with complex bodies push past a second (\texttt{rust-compression} 6064 ms, \texttt{rust-html-rewriter} 5104 ms, \texttt{regex} 5408 ms, \texttt{blind-sig} 3346 ms).

End-to-end wall-clock (\texttt{TtV}) is the metric that matters for cold-start workloads. LLVM JIT's TtV includes its instantiation latency in full. For short-run-time kernels, LLVM JIT's instantiation rivals or exceeds the benchmark execution time. \texttt{rust-html-rewriter} finishes in 5978 ms wall-clock, of which only 855 ms is benchmark work; the remaining 86\% is LLVM compile latency. \texttt{pulldown-cmark} and \texttt{rust-json} each spend roughly half their LLVM-JIT wall-clock time inside compilation. The tiered architecture's design goal is to reduce TtV for these workloads without losing steady-state throughput on the long-running kernels.

\section{Step 1: Promote Hot Callees}
\label{sec:eval:step1}

The earliest tier-2 design promotes only the function whose counter saturated, together with its direct callees. The batch's root is the threshold-tripper itself; the breadth-first traversal extends one hop to its direct callees. There is no walk up the call graph.

Across the thirty-one analyzed kernels, the geometric mean of \texttt{t1/t2} is \textbf{0.97$\times$} and the geometric mean of \texttt{LLVM/t2} is \textbf{0.79$\times$}. Tier-2 in this configuration is on average slightly slower than tier-1, and recovers approximately 79\% of LLVM JIT throughput.

The aggregate hides two distinct failure modes.

\paragraph{Hot-leaf failure.} For kernels in the hot-leaf category, the threshold-tripper is a small helper. Compiling only the leaf yields a tier-2 body whose call sites all dispatch back into tier-1 through the ABI bridge thunks introduced in Section~\ref{sec:design:context:thunks}. The optimizer cannot inline across the thunk boundary, so the calls remain expensive. The kernels in this category show the largest gaps to LLVM JIT: \texttt{shootout-ed25519} at 0.39$\times$, \texttt{blind-sig} at 0.33$\times$, \texttt{shootout-ctype} at 0.31$\times$. Each is below half of LLVM JIT throughput.

\paragraph{One-shot-caller failure.} For kernels in the one-shot-caller category, the entry-table swap fires after the outer function has already started running. The in-flight invocation runs the entire loop in tier-1 and returns. The swap is observable only on subsequent calls, and those calls never happen. \texttt{shootout-ctype}, \texttt{blind-sig}, \texttt{shootout-random}, and \texttt{shootout-base64} all sit well below the geometric mean for this reason. Several kernels in this category also regress relative to tier-1 (\texttt{shootout-ed25519} at 0.64$\times$ \texttt{t1/t2}, \texttt{shootout-ctype} at 0.51$\times$, \texttt{hashset} at 0.72$\times$), because the worker's compile cost is paid without a sufficient post-swap tail to amortize it.

The two improvements described in the following sections each address one of these failure modes.

\section{Step 2: Walk-Up and BFS Batch Composition}
\label{sec:eval:step2}

The improved batch composition walks up the static call graph from the threshold-tripper to find the hottest direct caller, then performs a bounded breadth-first traversal of that caller's direct callees. Callee inclusion is gated by the two-test admission rule described in Section~\ref{sec:design:tier2:batch}. The walk-up rolls the surrounding caller into the batch, so the optimizer has full bodies to inline into.

Across the thirty-one kernels, the geometric mean of \texttt{t1/t2} rises to \textbf{1.07$\times$} and the geometric mean of \texttt{LLVM/t2} rises to \textbf{0.88$\times$}. Tier-2 now beats tier-1 by approximately 7\% on average and recovers approximately 88\% of LLVM JIT throughput.

The hot-leaf kernels are the largest beneficiaries (Table~\ref{tbl:eval:walkup}). \texttt{shootout-ed25519} jumps from 0.39$\times$ to 0.69$\times$. \texttt{shootout-ctype} jumps from 0.31$\times$ to 0.64$\times$. \texttt{blind-sig} jumps from 0.33$\times$ to 0.50$\times$. The mechanism is the same in each case: the walk-up rolls the leaf's caller into the batch, the optimizer inlines short helpers that were previously trapped behind bridge thunks, and the resulting body runs closer to LLVM JIT speed.

\begin{table}[h]
\centering
\small
\begin{tabular}{l r r r}
\toprule
Kernel & Step 1 \texttt{LLVM/t2} & Step 2 \texttt{LLVM/t2} & Swing \\
\midrule
\texttt{shootout-ctype}    & 0.31$\times$ & 0.64$\times$ & $+0.33$ \\
\texttt{shootout-ed25519}  & 0.39$\times$ & 0.69$\times$ & $+0.30$ \\
\texttt{blind-sig}         & 0.33$\times$ & 0.50$\times$ & $+0.17$ \\
\texttt{rust-json}         & 0.63$\times$ & 0.79$\times$ & $+0.16$ \\
\texttt{rust-protobuf}     & 0.65$\times$ & 0.82$\times$ & $+0.17$ \\
\texttt{rust-compression}  & 0.78$\times$ & 0.89$\times$ & $+0.11$ \\
\bottomrule
\end{tabular}
\caption{Hot-leaf kernels with the largest \texttt{LLVM/t2} gain when the batch composition switches from leaf-only to walk-up plus BFS (Step 1 $\to$ Step 2).}
\label{tbl:eval:walkup}
\end{table}

The frequent-call kernels show a much smaller change. \texttt{shootout-keccak} moves from 1.00$\times$ to 1.00$\times$, \texttt{shootout-gimli} from 0.99$\times$ to 0.99$\times$, \texttt{blake3-scalar} from 0.97$\times$ to 0.99$\times$. These kernels were already in good shape at Step 1: function-entry promotion fires before the in-flight tail dominates the measurement, and the batch the optimizer sees is large enough to inline within.

The remaining gap kernels are one-shot callers. \texttt{shootout-ctype} reaches only 0.64$\times$, \texttt{blind-sig} only 0.50$\times$, \texttt{shootout-random} only 0.63$\times$. The function-entry promotion that walk-up and BFS produces is correct and fires; it simply has no effect on the in-flight invocation, and the in-flight invocation is the only invocation for these kernels. Closing this last group is the role of on-stack replacement.

\section{Step 3: On-Stack Replacement}
\label{sec:eval:step3}

The full configuration adds on-stack replacement (OSR) on top of walk-up and BFS. Tier-1 code now instruments each outermost-loop back-edge with the three-state diamond described in Section~\ref{sec:design:osr:diamond}. When the back-edge counter for a loop saturates, the worker thread compiles a continuation entry whose entry point is the loop header. The currently-executing tier-1 frame migrates into the continuation entry on the next iteration.

Across the thirty-one kernels, the geometric mean of \texttt{t1/t2} rises to \textbf{1.12$\times$} and the geometric mean of \texttt{LLVM/t2} rises to \textbf{0.92$\times$}. Tier-2 with OSR beats tier-1 by approximately 12\% on average and recovers approximately 92\% of LLVM JIT throughput.

The one-shot-caller kernels show the largest swings (Table~\ref{tbl:eval:osr}). The seven kernels with the biggest \texttt{LLVM/t2} jump from Step 2 to Step 3 are all from this category.

\begin{table}[h]
\centering
\small
\begin{tabular}{l r r r}
\toprule
Kernel & Step 2 \texttt{LLVM/t2} & Step 3 \texttt{LLVM/t2} & Swing \\
\midrule
\texttt{shootout-random}   & 0.63$\times$ & 1.00$\times$ & $+0.37$ \\
\texttt{shootout-ctype}    & 0.64$\times$ & 0.94$\times$ & $+0.30$ \\
\texttt{blind-sig}         & 0.50$\times$ & 0.76$\times$ & $+0.26$ \\
\texttt{shootout-matrix}   & 0.85$\times$ & 1.02$\times$ & $+0.17$ \\
\texttt{shootout-ratelimit}& 0.81$\times$ & 0.98$\times$ & $+0.17$ \\
\texttt{shootout-minicsv}  & 0.81$\times$ & 0.96$\times$ & $+0.15$ \\
\texttt{shootout-base64}   & 0.84$\times$ & 0.99$\times$ & $+0.15$ \\
\bottomrule
\end{tabular}
\caption{Kernels with the largest \texttt{LLVM/t2} gain when OSR is enabled (Step 2 $\to$ Step 3). All seven are from the one-shot-caller category.}
\label{tbl:eval:osr}
\end{table}

The mechanism behind each swing is identical. The outer function is called once. Almost all run-time is spent in one long-running loop. Without OSR, the function-entry slot swap fires after the loop has already started and is therefore observable only on a subsequent call that never happens. With OSR, the loop's back-edge counter saturates while the loop is still running, the worker compiles a continuation entry, and the tier-1 frame jumps into the continuation entry on the next iteration. The remainder of the loop runs in optimizing-tier code.

A small number of kernels regress when OSR is enabled. \texttt{rust-json} drops from 0.79$\times$ to 0.65$\times$. \texttt{rust-html-rewriter} drops from 0.71$\times$ to 0.64$\times$. \texttt{shootout-sieve} drops from 0.87$\times$ to 0.81$\times$. These kernels share two properties: the inner loop is short relative to the rest of the program, and the function is called frequently so that function-entry promotion was already capturing the hot path. The added back-edge instrumentation costs three operations per iteration on the post-saturation steady state, which is small in absolute terms but visible when the loop body itself is short. The OSR continuation compile also competes with regular tier-2 batches for worker time. The OSR-priority drain (Section~\ref{sec:design:tier2:worker}) ensures continuation entries land before the loop exits on the kernels for which OSR is the only mechanism that helps; the cost is a small delay for regular-tier-2 compiles on kernels for which OSR is unnecessary.

\section{Cold-Start and End-to-End Latency}
\label{sec:eval:tt}

The work-time analysis above is the steady-state metric. The cold-start metric is instantiation latency. The end-to-end metric is time-to-value.

Across the thirty-one kernels, the geometric mean of \texttt{Inst.Lat T2/T1} is \textbf{3.33$\times$} in the full configuration. The worker thread does more upfront work than tier-1 alone because it must allocate the OSR entry table, the back-edge counter array, and a per-thread frame buffer, but the bulk of optimization runs off the critical path. The ratio of LLVM JIT instantiation to tier-2 instantiation, \texttt{LLVM/T2 Inst.Lat}, has a geometric mean of \textbf{11.0$\times$} and ranges from approximately 3$\times$ on small modules to over 50$\times$ on the largest. LLVM JIT's per-instantiation cost is the structural cost that motivates the tiered architecture.

End-to-end \texttt{TtV} reflects the trade-off between compile cost and steady-state throughput. Across the thirty-one kernels, the geometric mean of \texttt{TtV LLVM/T2} is approximately \textbf{1.16$\times$}: the tiered architecture finishes the end-to-end workload approximately 16\% faster than whole-module LLVM JIT, even though its steady-state throughput is approximately 8\% lower. The gain comes from the saved instantiation latency.

The benefit is largest on short-run-time kernels. \texttt{rust-html-rewriter} finishes in 1806 ms under tier-2 versus 5978 ms under LLVM JIT, a 3.31$\times$ wall-clock win. \texttt{shootout-nestedloop} finishes 1.66$\times$ faster end-to-end. \texttt{regex} and \texttt{shootout-switch} both finish over 1.3$\times$ faster. These are exactly the workloads the cold-start architecture targets: request handlers that return in a second or two, where the full LLVM compile would otherwise dominate the user-visible response time.

The benefit shrinks on long-run-time kernels, because LLVM JIT's compile cost is amortized over more wasm execution. \texttt{shootout-ctype} \texttt{TtV} is approximately 1.00$\times$: tier-2 and LLVM JIT finish in essentially the same wall-clock time, by different paths. \texttt{shootout-sieve} \texttt{TtV} is 0.86$\times$, slightly slower end-to-end under tier-2 because the OSR steady-state regression on this kernel is not compensated by a large enough cold-start saving.


